\section{GlyphChain and Legal local completion and external release gates}
\label{sec:glyphchain-local-completion-external-gates}

This section records the final local implementation boundary.  The document-proof,
signing, delivery, consensus, recovery, template-governance and release-evidence
mechanisms required for a controlled GlyphChain Legal release are implemented and
covered by the retained regression baseline.  The remaining thirty-four tasks are
not unfinished synthetic-development exercises: each requires evidence from a live
provider, an independently operated system or reviewer, a qualified Legal reviewer,
a customer pilot, or an authorised production decision.

This distinction prevents Tessaris from manufacturing the evidence needed to
approve itself.  A local feature is complete when it can securely collect and
verify the external evidence.  The associated release gate remains open until that
evidence actually exists for the exact release candidate.

\subsection{Completed local product foundation}

The following capabilities are available to the Legal workspace and release
process:

\begin{itemize}
  \item encrypted Legal matters, documents, templates, envelopes, audit events and
        retention state, with production key custody separated from tracked data;
  \item controlled signature-envelope preparation, signer sequencing, access-token
        boundaries, evidence capture, completion certificates, signed-copy
        manifests and downloadable private evidence packages;
  \item portable \texttt{.tessaris-proof} receipts and a standalone local verifier
        for document hash, device Ed25519 signature, GlyphChain transaction,
        Merkle inclusion and trust-anchored quorum-finality evidence;
  \item a durable, idempotent proof-submission outbox with restart recovery,
        exponential retry, duplicate-chain lookup and workspace isolation;
  \item verified consensus restart recovery, peer identity handshakes, malformed
        message rejection, congestion controls, bounded catch-up and deterministic
        finality-certificate validation;
  \item Resend and Microsoft 365 delivery adapters with authenticated-event
        ingestion, exact-recipient state, controlled retry and a content-free
        delivery-qualification report;
  \item immutable trust-release construction, an independently pinnable checkpoint
        identifier and a bounded clean-device verification exercise;
  \item representative network, protocol-upgrade, relay/partition, validator-
        removal and sustained-soak qualification tools whose reports exclude Legal
        document content;
  \item encrypted application-account recovery kits, root-signed public account
        descriptors, monotonic recovery epochs, device revocation and application-
        session authority checks;
  \item exact-hash Legal template approval, jurisdiction binding, withdrawal,
        supersession and review-required fallback behaviour; and
  \item a signed exact-source baseline and fail-closed release-readiness evaluator
        covering every remaining production task.
\end{itemize}

These functions provide document integrity, attributable consent evidence,
controlled workflow and verifiable ledger assurance.  They do not by themselves
establish universal legal enforceability, qualified-electronic-signature status,
professional Legal advice, independent network operation or production approval.

\subsection{Canonical external task register}

The written remaining-work checklist and the machine release policy share the same
stable identifiers.  A regression extracts the identifiers from the checklist and
requires exact equality with the runtime register; duplicate, omitted and unknown
identifiers fail the test.  The register is:

\begin{center}
\begin{tabular}{lll}
\hline
\textbf{Identifiers} & \textbf{External workstream} & \textbf{Count} \\
\hline
\texttt{LD-01--LD-05} & Live signing delivery & 5 \\
\texttt{TR-01--TR-03} & Independent trust publication & 3 \\
\texttt{NET-01--NET-05} & Distributed-network production validation & 5 \\
\texttt{SEC-01--SEC-04} & Production security and recovery & 4 \\
\texttt{LEGAL-01--LEGAL-07} & Qualified Legal content & 7 \\
\texttt{REL-01--REL-10} & Controlled deployment and release & 10 \\
\hline
& \textbf{Total} & \textbf{34} \\
\hline
\end{tabular}
\end{center}

\subsection{Live delivery closure evidence}

Tasks \texttt{LD-01} to \texttt{LD-04} require authenticated events from the
production Resend account and Microsoft 365 tenant.  The local implementation can
verify Resend webhook signatures, ingest Microsoft Message Trace evidence, preserve
the exact Internet Message-ID and recipient, distinguish delivery, delay and
failure, and require renewed authority before retry.  It deliberately does not
promote provider submission to delivery or infer a detailed bounce reason from a
coarse failure state.

Task \texttt{LD-05} closes only when the production-mode qualification report is
retained for the exact release and signed by a delivery operator and a distinct
independent security reviewer.  Synthetic callback fixtures prove parser and state
machine behaviour but cannot replace live provider evidence.

\subsection{Independent trust and clean-device evidence}

Task \texttt{TR-01} requires the immutable checkpoint, offline verifier and
digest-bound manifest to be published through the approved versioned release
channel.  Task \texttt{TR-02} requires the checkpoint identifier to be published
through a genuinely independent channel.  The release package cannot appoint its
own trust anchor merely by carrying both checkpoint and pin.

Task \texttt{TR-03} requires a separate clean physical or virtual device to verify
a real private Legal evidence package with an independently obtained pin.  The
implemented exercise retains only content-free digests and outcomes; the private
document remains with the authorised verifier.

\subsection{Distributed-network and security evidence}

The local qualification runners model representative node profiles, reconnects,
relay outages, partitions, protocol activation, validator removal and sustained
load.  Tasks \texttt{NET-01} to \texttt{NET-05} nevertheless require independent
operators, real network boundaries, a production-like public relay and an
independently observed proof that no Legal document content crosses the relay.
Several logical nodes on one development Mac are regression evidence, not a claim
of production decentralisation.

Task \texttt{SEC-01} requires the bounded current-account-descriptor registry to be
deployed under separate operational control and used by every production session
issuer.  A recovered account must advance its epoch, authorise only the replacement
device and cause stale epochs and removed device identifiers to fail closed.
Validator-finality keys remain in deployment secret custody and outside the account
recovery kit.

Tasks \texttt{SEC-02} and \texttt{SEC-03} require independent application-security
and consensus/cryptography review, including resolution and retest of findings.
Task \texttt{SEC-04} requires the passing bounded soak to be repeated across
representative independently operated hosts, real restarts and intermittent
networks under production-like load.

\subsection{Qualified Legal content boundary}

Tasks \texttt{LEGAL-01} to \texttt{LEGAL-06} require qualified reviewers and
approved wording for every enabled jurisdiction and template family: customer,
supplier, contractor and employment agreements, terms and conditions, privacy
notices and data-processing agreements.  The product cannot appoint a qualified
reviewer or approve its own Legal wording.

Task \texttt{LEGAL-07} uses the implemented encrypted template library to bind the
exact reviewed bytes to reviewer identity, jurisdiction, approval date, SHA-256
content hash and supersession rules.  A changed byte creates a different version
and requires review.  Withdrawn, superseded, unmatched or unapproved templates
fall back to a clearly labelled review-required draft; they cannot silently power
automated issue or signature.

\subsection{Controlled release and measured pilot}

Tasks \texttt{REL-01} to \texttt{REL-03} require publication of the versioned
assurance, privacy, participation, recovery, operator and electronic-signature
claim boundaries.  Product wording must continue to distinguish local evidence,
ledger commitment, internally valid but unanchored signatures, and independently
trust-anchored quorum finality.  It must not claim qualified-electronic-signature
status unless that separate regulatory route is completed or an approved trust
service supplies it.

Tasks \texttt{REL-04} to \texttt{REL-07} require an internal multi-device
deployment, an approved closed customer pilot, measured proof-finality, delivery,
verification, restore and recovery outcomes, and resolution plus retest of live
security, reliability and usability findings.  Tasks \texttt{REL-08} and
\texttt{REL-09} then freeze the exact release, independently pin its trust material
and collect distinct final security, Legal product and release approvals.

\subsection{Exact-release authorization algorithm}

Each external evidence body includes the task identifier $t$, production
environment, observation time, result, statement, typed artifact digests and the
exact release-manifest digest $h_M$.  Required authorities sign a domain-separated
canonical body with Ed25519.  Trust roots define authorised roles and whether a
signer is independent.

For task $t$, let $B_t$, $A_t$, $R_t$, $I_t$ and $S_t$ denote valid release
binding, artifact evidence, role signatures, required independence and separation
of duties respectively.  The task result is

\begin{equation}
  G_t = B_t \land A_t \land R_t \land I_t \land S_t.
\end{equation}

The production decision for the fixed register $T$ is

\begin{equation}
  G_{\mathrm{production}} =
  \left(\bigwedge_{t \in T} G_t\right)
  \land (U=\varnothing)
  \land (D=\varnothing),
\end{equation}

where $U$ and $D$ are the unknown and duplicate identifier sets.  Multi-role tasks
require distinct signers, and an independent task requires an independent signer
who actually fulfils one of its required roles.  There is no percentage threshold,
implicit waiver, development override or self-attested substitute.

Task \texttt{REL-10} is satisfied only after this algorithm reports all thirty-four
tasks passed for the frozen release.  The readiness report is written with
owner-only permissions, refuses overwrite and symbolic-link destinations, and
contains a SHA-256 commitment over all decision fields.

\subsection{Current release state and permitted claim}

The production evidence and trust-root inventories intentionally contain no
fabricated authorities or events.  The current exact-release evaluation therefore
returns \texttt{BLOCKED\_EXTERNAL\_GATES}, zero passed, thirty-four unmet and
\texttt{production\_release\_authorized=false}.  This is the expected safe result
after local implementation completion.

The permitted statement is: ``The local GlyphChain and Legal protocol foundation,
evidence collectors and release-verification mechanisms are implemented and pass
the retained regression baseline; production remains gated on live, independent,
qualified and pilot evidence for the exact release.''  Tessaris must not shorten
this to ``blockchain certified'', ``legally certified'', ``qualified signature''
or ``production ready'' until the corresponding external gates have genuinely
closed.

\subsection{Release artefacts and operator sequence}

The authoritative local artefacts are the remaining-work checklist, this local-
completion section, the signed baseline report, the empty evidence and trust-root
inventory templates, the versioned readiness report, and the standalone assurance,
recovery, delivery, trust-publication, network and security sections.

For each production candidate, operators shall:

\begin{enumerate}
  \item run the full retained regression suite and refuse a failing baseline;
  \item generate and sign the exact source/release manifest;
  \item collect external evidence without private Legal document content;
  \item verify artifact digests and role-scoped attestations against independently
        administered trust roots;
  \item run the readiness evaluator for the exact manifest;
  \item resolve every unmet, unknown or duplicate task; and
  \item publish or deploy only when all thirty-four tasks pass and the final report
        explicitly authorises the candidate.
\end{enumerate}

This sequence completes all work that can safely be automated locally while
preserving the human, professional and independent controls that give the release
evidence its meaning.
